% preamble.tex

% Türkçe karakter desteği ve dil ayarları
\usepackage[utf8]{inputenc}
\usepackage[T1]{fontenc}
\usepackage[turkish]{babel}
% babel Türkçe'de eşittir işaretinin shorthand olmasını sevmiyor gibi
\AtBeginDocument{\shorthandoff{=}}

% Satır aralığı: 1.5 (Türkpatent zorunluluğu)
\usepackage{setspace}
\onehalfspacing

% Metin sayfaları için sayfa kenar boşlukları (Üst 2-4 cm, Sol 2.5-4 cm, Sağ 2-3 cm, Alt 2 cm)
\usepackage[top=3cm, left=3cm, right=2.5cm, bottom=2cm]{geometry}

% Satır numaralandırması (Her 5 satırda bir)
\usepackage[left]{lineno}
\modulolinenumbers[4]

% Sayfa numaralandırması (Alt orta kısımda)
\usepackage{fancyhdr}
\pagestyle{fancy}
\fancyhf{} % Üst ve alt bilgileri temizle
\renewcommand{\headrulewidth}{0pt} % Üst çizgi kalınlığını sıfırla
\cfoot{\thepage} % Sayfa numarasını alta ortala 

% Resim ekleme desteği ve resimlerin aranacağı klasör
\usepackage{graphicx}
\graphicspath{{sekiller/}}

% Resim sayfaları (04-resimler.tex) için özel kenar boşluğu ve sayfa numaralandırması
% ayarlarını ana belgede veya ilgili dosyanın içinde yöneteceğiz.
